%% resume.tex — Mathias Dick
%% Paste into https://overleaf.com (new project → paste code)
%% Requires: fontawesome5, geometry, titlesec, enumitem, hyperref, xcolor, multicol

\documentclass[10pt, a4paper]{article}

% ─── Packages ────────────────────────────────────────────────────────────────
\usepackage[a4paper, top=1.6cm, bottom=1.6cm, left=1.8cm, right=1.8cm]{geometry}
\usepackage{fontawesome5}
\usepackage{titlesec}
\usepackage{enumitem}
\usepackage{hyperref}
\usepackage{xcolor}
\usepackage{multicol}
\usepackage{paracol}
\usepackage{parskip}
\usepackage{microtype}
\usepackage{cmap}
\usepackage[T1]{fontenc}
\usepackage[default]{sourcesanspro} % Clean, modern sans-serif font
\usepackage{graphicx}

% ─── Colors ──────────────────────────────────────────────────────────────────
\definecolor{accent}{HTML}{2563EB}   % blue accent
\definecolor{subtle}{HTML}{6B7280}   % grey subtitles

% ─── Hyperlinks ──────────────────────────────────────────────────────────────
\hypersetup{colorlinks=true, urlcolor=accent, hidelinks}

% ─── Remove page numbering ───────────────────────────────────────────────────
\pagestyle{empty}

% ─── Disable Hyphenation ─────────────────────────────────────────────────────
\hyphenpenalty=10000
\exhyphenpenalty=10000

% ─── Section style ───────────────────────────────────────────────────────────
\titleformat{\section}
  {\large\bfseries\color{accent}}
  {}{0em}
  {}
  [\vspace{-0.4em}\rule{\linewidth}{0.5pt}\vspace{0.2em}]

\titlespacing*{\section}{0pt}{0.8em}{0.4em}

% ─── List style ──────────────────────────────────────────────────────────────
\setlist[itemize]{
  leftmargin=1.2em,
  itemsep=0.1em,
  parsep=0em,
  topsep=0.2em
}

% ─── Custom commands ─────────────────────────────────────────────────────────

% \entry{Title}{Subtitle / Period}{Description or items}
\newcommand{\entry}[3]{%
  \textbf{\large #1} \hfill \textcolor{subtle}{\small #2}%
  \def\temp{#3}\ifx\temp\empty\else\\[0.1em]\textbf{#3}\fi
}

% \project{Title}{Tools}
\newcommand{\project}[2]{%
  \textbf{\textit{#1}} \hfill \textcolor{subtle}{\footnotesize \textbf{Tech:} #2}\vspace{0.1em}
}

% ─── Ragged Right for Skills ─────────────────────────────────────────────────
\newcommand{\skill}[2]{%
  \textbf{#1}\\[0.15em]
  {\raggedright #2\par}
}

% ─────────────────────────────────────────────────────────────────────────────
\begin{document}

% ═══ HEADER + SUMMARY ═══════════════════════════════════════════════════════
% Place your photo as photo.jpg in the same folder as this .tex file
\begin{minipage}[t]{0.76\linewidth}
  \vspace{0pt}
  {\Huge\bfseries Mathias Dick}\\[0.4em]
  \small
  \faMapMarker*\enspace Rastatt, Germany\quad
  \faEnvelope\enspace \href{mailto:mathias.dick16@gmail.com}{mathias.dick16@gmail.com}\\[0.3em]
  \faGlobe\enspace Portfolio: \href{https://mathiasd.de?utm_source=cv}{mathiasd.de}\quad
  \faGithub\enspace \href{https://github.com/MathiasDick}{github.com/MathiasDick}\quad
  \faLinkedin\enspace \href{https://www.linkedin.com/in/mathias-dick/}{linkedin.com/in/mathias-dick}
  \vspace{0.5em}

  {\large\bfseries\color{accent} Zusammenfassung}\\[-0.3em]
  \rule{\linewidth}{0.5pt}\\[0.3em]
  \small
  Multidisziplinärer Ingenieur, der Softwareentwicklung, Elektronik und praktische Mechanik-Kenntnisse vereint.
  Derzeit Dualer Student bei der \textbf{Mercedes-Benz AG}, mit Fokus auf der Verbindung von Full-Stack- und Desktop-Anwendungsentwicklung mit Automobiltechnik.
  Starker Fokus auf die Entwicklung zuverlässiger, produktionsreifer Tools — von KI-gestützter Fehlerticket-Generierung bis hin zu industriellen Netzwerk-Monitoring-Systemen.\\[0.4em]
  \href{https://mathiasd.de/work/?utm_source=cv_summary}{\textit{\footnotesize\textcolor{subtle}{\faLink\enspace Detaillierte technische Berichte, Architekturdiagramme und Fotos meiner Projekte sind unter mathiasd.de/work verfügbar}}}
\end{minipage}
\hfill
\begin{minipage}[t]{0.20\linewidth}
  \vspace{0pt}
  \includegraphics[width=\linewidth, keepaspectratio=true]{photo.jpg}
\end{minipage}
\vspace{0.6em}

% ═══ TWO-COLUMN BODY ═════════════════════════════════════════════════════════
% \setlength{\columnsep}{1.5em}
% \columnratio{0.64}
% \begin{paracol}{2}
\begin{minipage}[t]{0.64\linewidth}

% ─── LEFT COLUMN ─────────────────────────────────────────────────────────────

\section{Berufserfahrung}

\entry{Dualer Student}{Okt 2023 – Sep 2026 (Voraussichtlich)}{Mercedes-Benz AG · Rastatt, Deutschland}\\[0.2em]
\project{KI-Ticket-Assistent}{Python · React · Flask · GPT-4o · Whisper · RAG}
\begin{itemize}
  \item Entwicklung einer hybriden KI-Pipeline mit lokaler Whisper-Transkription und GPT-4o zur Automatisierung der formalen Fehlerticket-Generierung, wodurch die manuelle Eingabezeit für über 80 Tickets wöchentlich um ca. 50\% reduziert wurde.
  \item Implementierung einer RAG-basierten Duplikaterkennung über eine vektorisierte Fehlerhistorie, um redundante Einreichungen zu verhindern.
  \item Bereitstellung als eigenständige ausführbare Datei, die importfertige Tickets und HTML-Berichte generiert.
\end{itemize}

\vspace{0.4em}

\project{Industrieller Redundanz-Monitor}{Python · AsyncIO · Flask · React · SNMP · SQLite}
\begin{itemize}
  \item Entwicklung eines Monitoring-Systems zur Erkennung stiller Redundanzausfälle über 500+ Industrie-Switches hinweg, um Ausfallzeiten zu vermeiden.
  \item Entwicklung eines asynchronen SNMP-Scanners mit Auto-Discovery, der die manuelle Bestandsverwaltung überflüssig macht.
  \item Entwicklung eines virtualisierten React-Frontends und einer sicheren, JWT-geschützten API, bereitgestellt über IIS.
\end{itemize}

\vspace{0.4em}

\entry{Freiberuflicher Entwickler}{Jan 2024 – Heute}{}\\[0.2em]
\project{Markt-Monitor}{Python · PostgreSQL · SQLAlchemy · n8n · Linux}
\begin{itemize}
  \item Entwicklung einer Produktions-Datenpipeline, die über 170.000 Produkte auf großen E-Commerce-Plattformen scannt, um Trendanalysen und marktübergreifende Preisvergleiche zu ermöglichen.
  \item Implementierung nebenläufiger Worker mit Prozess-Locking und KI-gesteuerter Datenanreicherung über n8n, wodurch über 50.000 falsch gelistete Artikel identifiziert wurden.
  \item Selbst gehostet auf Linux, mit eigenständiger Verwaltung von Deployments, Updates und Workflow-Monitoring.
\end{itemize}

\vspace{0.4em}

\section{Ausgewählte Projekte}

\textbf{Audi Motorinstandsetzung} \hfill \textcolor{subtle}{\small Jun 2025}\\[0.15em]
\textcolor{subtle}{\footnotesize Diagnose · Präzisionsmesstechnik}\\
\textit{\footnotesize Diagnose und Reparatur eines katastrophalen Ventilschadens an einem 2.0 TFSI Motor. Anwendung systematischer Fehlersuche und Werkstoleranzen für eine zuverlässige Instandsetzung.}

\vspace{0.5em}

\textbf{RFID Smart-Türschloss} \hfill \textcolor{subtle}{\small Feb 2024}\\[0.15em]
\textcolor{subtle}{\footnotesize Arduino · C++ · 3D-Druck · CAD}\\
\textit{\footnotesize Entwicklung eines schlüssellosen Zugangssystems mit RFID und Schrittmotorsteuerung, inklusive maßgeschneiderter, zerstörungsfrei montierbarer 3D-Druck-Halterungen.}

% ─── RIGHT COLUMN ────────────────────────────────────────────────────────────
\end{minipage}
\hfill
\begin{minipage}[t]{0.32\linewidth}

\section{Ausbildung}

\textbf{B.Eng. Embedded Systems}\\
\textcolor{subtle}{\small Studienrichtung Fahrzeugelektronik}\\
\textcolor{subtle}{\small DHBW Ravensburg}\\
\textcolor{subtle}{\small Campus Friedrichshafen}\\
\textcolor{subtle}{\small Okt 2023 – Sep 2026 (Voraussichtlich)}

\vspace{0.5em}

\textit{\footnotesize\textcolor{subtle}{Duales Studium in Kooperation mit der Mercedes-Benz AG — Wechsel zwischen Praxisphasen im Unternehmen und Theoriephasen an der Hochschule.}}

\vspace{0.5em}

\section{Fähigkeiten}

\skill{Sprachen \& Runtimes}{Python 3.10+, JavaScript / TypeScript, C++ (Embedded), SQL}

\vspace{0.4em}
\skill{Frameworks \& Bibliotheken}{React, Redux Toolkit, Vite, Flask, Waitress, SQLAlchemy, Alembic, Pandas, AsyncIO, PySide6, Selenium}

\vspace{0.4em}
\skill{KI / ML}{Azure OpenAI (GPT-4o), Whisper (lokal), RAG / Vektor-Embeddings, Prompt Engineering}

\vspace{0.4em}
\skill{Infrastruktur}{PostgreSQL, SQLite, n8n, IIS, Windows Server, Linux, Git, PyInstaller}

\vspace{0.4em}
\skill{Hardware \& Werkzeuge}{Arduino / Embedded Systems, SNMP, Fusion 360, Elektronikreparatur \& Diagnose}

\section{Sprachen}

\begin{itemize}[leftmargin=0.8em]
  \item \textbf{Deutsch} — Muttersprache
  \item \textbf{Englisch} — Fließend
  \item \textbf{Russisch} — Gute Kenntnisse
  \item \textbf{Japanisch} — Mittelstufe (JLPT N2)
\end{itemize}

\end{minipage}

\end{document}
